\chapter{Arquitectura}
\section*{2a Arquitectura web}
\subsection*{Dudas de conocimiento}
Mi mayor problema en este apartado más que entender la información y los conceptos es saber que es lo importante y la información que do obtener de esta.


Las ideas que hay que sacar de esta parte son:
\begin{itemize}
    \item  Diferencia sitios estaticos vs sitios dinamicos y como funciona cada uno 
    \item Principios diseño software: \begin{itemize}
        \item SOLID
        \item KISS (KEEP IT SUPER SIMPLE)
        \item YAGNI (YOU AREN´T GONNA NEED IT)
        \item DRY (DON'T REPEAT YOURSELF)
    \end{itemize}
    \item  Razones por las que reutilizar codigo. 
\end{itemize}
\subsection*{Dudas de contenido para el examen}
\begin{enumerate}
    \item ¿Hay que saber los ejemplos de CGI que se muestran en las diapos de la 28 a la 32?
    \item Los principios de diseño software ¿hay que saberlos de memoria o solo entenderlos y saber aplicarlos?
\end{enumerate}



\section*{2b Arquitectura web}
\subsection*{Dudas de conocimiento}
\begin{itemize}
    \item ¿Los filters y los controllers son de a capa de presentación?
    \item No entiendo el problema de la diapositiva 25,26 de Model 1, me refiero, si se produce un error en el proceso no es su responsabilidad, no? simplemente se reenvia a una pagina de error, me refieor, si ya has enviado la mitad de la info y se produce un error con la base de datos no es responsabilidad de la capa de presentación, simplenente redirije a la pagina de error.
    \item No entiendo muy bien las diferencias de la diapositiva 34 de diferencias entre Model 1 y 2, del data push vs Data pull. Lo que he entendido es que el model 1 pide los datos, es decir hace una llamada a Persistencia y el model 2 us el controlador para enviarlos directamente sin que los pida porque los necesita. 
    \item En resumen para entenderlo, el model 2 la diferencia pricipal es el controlador que hace de capa intermedia y reparte los recursos y esto le da varias ventajas como no introducir tanto codigo, no repetir codigo o no empezar a enviar una vista y que luego de error. 
\end{itemize}

\textbf{No es una duda, pero hay una errata en la diapositiva 40,, escribe la mantenibilidad junto en vez de separado}

Las ideas que hay que sacar de esta parte son:
\begin{itemize}
    \item  Model 0
    \item Model 1
    \item Patrón MVC (Model 2)
\end{itemize}

\subsection*{Dudas de contenido para el examen}

\begin{enumerate}
    \item ¿Hay que saberse de memoria los distintos Frameworks de MVC que se muestran en la diapo 48?
\end{enumerate}

\section*{2c Arquitectura web- Tecnologias Java}
\subsection*{Dudas de conocimiento}

Las ideas que hay que sacar de esta parte son:
\begin{itemize}
    \item  Entorno Java EE: aplicación web, saberse que es y que contiene y como se puede desplegar.
    \item Apache tomcat 
    \item Integración del servidor
    \item Concepto de Servlet
    \item Concepto de Context
    \item Archivo WAR concepto y funcionamiento 
    \item dir WEB-INF y su contenido
    \item web.xml y su contenido
\end{itemize}

\subsection*{Dudas de contenido para el examen}
\begin{itemize}
    \item De la diapo 4 hay que saberse los cambios que se han producidp de una version a otra de java ee y jakarta?
    \item Lo mismo sobre las API Servlet de diapo 14. 
    \item Hay que saberse de memoria la estructura de una archivo war como la diapo 18 o solo entenderla y saber que es cada cosa y poder hubicarla donde debe ir?
\end{itemize}


\section*{2d Estilos Arquitectonicos}
\subsection*{Dudas de conocimiento}
\begin{itemize}
    \item No entiendo las diferencias entre frameworks y arquitecturas, lo que he entendido es que las arquitecturas son un conjunto de patrones y principios de diseño que se pueden aplicar a cualquier proyecto y los frameworks son una implementación concreta de una arquitectura, es decir, un conjunto de herramientas y librerías que facilitan el desarrollo de aplicaciones siguiendo una arquitectura específica, pero no estoy seguro si es eso a lo que se refiere. 
\end{itemize}
Las ideas que hay que sacar de esta parte son:
\begin{itemize}
    \item Frameworks
    \item Big Ball of Mud (BBoM)
\end{itemize}
\subsection*{Dudas de contenido para el examen}
\begin{itemize}
    \item ¿Hay que saberse de memoria los nombre e infor de la diapo 5 a la 10 de las distintas estructuras?
end{itemize}
